% ======================================================================
% TUM-Farben
% ======================================================================
% TUM blue
\definecolor{tum blue}{HTML}{0065BD}
% Additional blue color tones
% 1 is the lightest (light blue)
% 5 is the darkest (deep navy)
\definecolor{tum blue 1}{HTML}{98C6EA}
\definecolor{tum blue 2}{HTML}{64A0C8}
\definecolor{tum blue 3}{HTML}{0073CF}
\definecolor{tum blue 4}{HTML}{005293}
\definecolor{tum blue 5}{HTML}{003359}

% TUM accent colors
\definecolor{tum green}{HTML}{A2AD00}
\definecolor{tum orange}{HTML}{E37222}
\definecolor{tum ivory}{HTML}{DAD7CB}

% TUM diagram colors
\definecolor{tum dia violet}{HTML}{69085A}
\definecolor{tum dia dark blue}{HTML}{0F1B5F}
\definecolor{tum dia turquoise}{HTML}{00778A}
\definecolor{tum dia dark green}{HTML}{007C30}
\definecolor{tum dia light green}{HTML}{679A1D}
\definecolor{tum dia light yellow}{HTML}{FFDC00}
\definecolor{tum dia dark yellow}{HTML}{F9BA00}
\definecolor{tum dia dark orange}{HTML}{D64C13}
\definecolor{tum dia red}{HTML}{C4071B}
\definecolor{tum dia dark red}{HTML}{9C0D16}
% ======================================================================


% ======================================================================
% Überschrift
% ======================================================================
\newcommand{\Ueberschrift}{%
  {\Large\bf%
    \noindent\TypDerArbeit:\\[5 pt]%
    \noindent\TitelDerArbeit}

  \vspace*{14 pt}
}
% ======================================================================


% ======================================================================
% Einrückungen ein/aus
% ======================================================================
\newcommand{\EinrueckungEin}{\setlength\parindent{1.2 em}}
\newcommand{\EinrueckungAus}{\setlength\parindent{0 em}}
% ======================================================================



% ======================================================================
% Abschnitte ohne Nummerierung
% ======================================================================
\let\sectionold\section
\let\subsectionold\subsection
\let\subsubsectionold\subsubsection

\renewcommand{\section}[2][\doSection]{%
  \def\doSection{#2}
  \phantomsection%
  \sectionold*{#2}%
  \addcontentsline{toc}{section}{#1}%
}
\renewcommand{\subsection}[2][\doSubsection]{%
  \def\doSubsection{#2}
  \phantomsection%
  \subsectionold*{#2}%
  \addcontentsline{toc}{subsection}{#1}%
}
\renewcommand{\subsubsection}[2][\doSubsubsection]{%
  \def\doSubsubsection{#2}
  \phantomsection%
  \subsubsectionold*{#2}%
  \addcontentsline{toc}{subsubsection}{#1}%
}
% ======================================================================

% ======================================================================
% Eigene Befehle
% ======================================================================
% 1) Symbols 
\newcommand{\Reals}{\mathbb{R}}
\newcommand{\Complex}{\mathbb{C}}
\newcommand{\Htwo}{\mathcal{H}_2}
\newcommand{\Hinf}{\mathcal{H}_\infty}
\newcommand{\Ltwo}{\mathcal{L}_2}
\newcommand{\Linf}{\mathcal{L}_\infty}
\renewcommand{\th}{^{\text{th}}}
\newcommand{\BigO}[1]{\mathcal{O}\left(#1\right)}

\newcommand{\defeq}{\vcentcolon=}
\newcommand{\logicand}{\wedge}
\newcommand{\logicor}{\vee}

% 2) Matrix operations
\newcommand{\trans}[1]{#1^\top}
\newcommand{\htrans}[1]{#1^H}
\newcommand{\hermit}{^H}
\newcommand{\inv}[1]{#1^{-1}}
\newcommand{\invt}[1]{#1^{-\top}}
\newcommand{\pinv}[1]{#1^{\dag}}
\newcommand{\pd}{\succ}
\newcommand{\nd}{\prec}



% 3) Operators
\newcommand\abs[1]{\left|#1\right|}
\DeclareMathOperator{\vecspan}{span}
\newcommand\vspan[1]{\vecspan\left(#1\right)}
\newcommand\vecspace[1]{\mathcal{#1}}
\DeclareMathOperator{\diagmat}{diag}
\newcommand\diag[1]{\diagmat\left(#1\right)}
\DeclareMathOperator{\matrixrank}{rank}
\newcommand\rank[1]{\matrixrank #1}
\DeclareMathOperator{\matrixrang}{rang}
\newcommand\rang[1]{\matrixrang\left(#1\right)}

\newcommand{\norm}[1]{\left\lVert#1\right\rVert}
\newcommand{\diff}[2]{\frac{\mathrm{d}\, #1}{\mathrm{d}\, #2}}
\newcommand{\innProd}[2]{\left\langle#1, \; #2\right\rangle}
\newcommand{\range}[1]{\mathcal{R}\left(#1\right)}

% 4) Formatting
\newcommand{\ts}{\!} %negative space for formulas in text
\newcommand{\nix}{\mbox{ }} %empty space



%---- MOR comamnds
\newcommand{\krylov}[3]{\vecspace{K}_{#1}\left(#2, #3\right)}
%\newcommand{\shift}{\sigma}
\newcommand{\As}{A_{\shift}}
\newcommand{\Asx}{A-\shift E}
\newcommand{\inputkrylov}{\krylov{n}{\inv{\As} E}{\inv{\As}B}}
\newcommand{\outputkrylov}{\krylov{n}{\invt{\As} \trans{E}}{\invt{\As} \trans{C}}}
\newcommand{\moment}{C\left(\inv{\As}E\right)^i \inv{\As} B}

\newcommand{\pencil}{\mu E - A}

\newcommand{\GramC}{\mathcal{W}_c}
\newcommand{\GramO}{\mathcal{W}_o}

\newcommand{\fo}{n}
\newcommand{\ro}{q}
% ======================================================================
